% !TEX program = xelatex
% !TEX options = -synctex=1 -interaction=nonstopmode -file-line-error
% example-figures.tex
% Standalone demonstration of rossfigures: figurepanel with \panel,
% sidebyside, fullwidthfigure, and figurenote.

\documentclass[11pt, a4paper]{article}
\usepackage[T1]{fontenc}
\usepackage{mathpazo}
\usepackage[margin=25mm]{geometry}
\usepackage{rosspalettes}   % for themed muted colour
\usepackage{rossfigures}

\title{\texttt{rossfigures} --- Feature Demonstration}
\author{Ross Dunne}
\date{2026}

\begin{document}
\maketitle

\section*{1. figurepanel environment}

The \texttt{figurepanel} environment groups sub-figures with an overall
caption. Each \verb|\panel| places a labelled sub-figure at the
specified width (default 0.48\verb|\textwidth|).

\begin{figurepanel}[RBANS index score distributions by diagnostic group.]
  \panel{example-image-a}{(A) Immediate Memory}
  \panel{example-image-b}{(B) Delayed Memory}
\end{figurepanel}
\figurenote{Box plots show median, IQR, and 1.5$\times$IQR whiskers.
  Outliers plotted individually. HC = healthy control; VaD = vascular
  dementia.}

\section*{2. sidebyside}

The \verb|\sidebyside| command places two captioned images side by side
without needing the panel environment.

\sidebyside
  {example-image-a}{Scatter plot of SWS\% vs.\ delayed memory.}
  {example-image-b}{Mediation path diagram with standardised coefficients.}
\figurenote{Dashed lines indicate non-significant paths ($p > .05$).}

\section*{3. fullwidthfigure}

A single image spanning the full text width with a caption.

\fullwidthfigure{example-image}{%
  White-matter lesion probability map overlaid on the MNI-152 template.
  Colour scale represents lesion frequency across the VaD sample
  ($n = 110$).}
\figurenote{Axial slices at $z = -10$, 0, 10, 20, 30, and 40\,mm.
  Images are shown in radiological convention (left = right).}
\source{Data from the Limerick Vascular Ageing Study (2024).}

\section*{4. figuregrid}

For three or more panels use \verb|\figuregrid{cols}{...}| to
compute cell widths automatically.

\figuregrid{3}{%
  \panel[\rf@cellwidth]{example-image-a}{(A) Attention}
  \panel[\rf@cellwidth]{example-image-b}{(B) Language}
  \panel[\rf@cellwidth]{example-image}{(C) Visuospatial}
}
\figurenote{Error bars represent 95\% confidence intervals.}

\end{document}
